\documentclass[usenames,dvipsnames,notes]{beamer}
%\documentclass[usenames,dvipsnames,notes=only]{beamer}
%\documentclass[usenames,dvipsnames]{beamer}
% Use one of the three above, to keep or remove or get notes.

\usepackage{graphicx,hyperref,url}
\usepackage{xcolor}
\usepackage[absolute,overlay]{textpos}
\usepackage{ucs}
%\usepackage[utf8]{inputenc}
\usepackage{fontspec}
\usepackage{xunicode} %% loading this first to avoid clash with bidi/arabic

%%% For language switching -- like babel, but for xelatex
\usepackage{polyglossia}
\setmainlanguage{english}
\setotherlanguages{hindi} %% or other languages
\newfontfamily\hindifont{Noto Sans Devanagari}[Script=Devanagari]

\usepackage{amsmath}
\usepackage{amsfonts}
\usepackage{amssymb}
\usepackage{amsthm}
\usepackage{caption}
\usepackage{multicol}
\usepackage{subfigure}
\usepackage{graphicx}
\usepackage{transparent}
\usepackage{natbib}
\usepackage{pstricks}
\usepackage{bbold}
\setbeamertemplate{navigation symbols}{}
\newcommand*{\vcenteredhbox}[1]{\begin{tabular}{@{}c@{}}#1\end{tabular}}
%%% : End of RSJ CS BITS GOA Themes
%% Baskar's hack for hiding slide numbers to avoid spooking deans on number of slides:
\newcommand\numbered{\setbeamertemplate{footline}{%
		\raisebox{5pt}{\makebox[\paperwidth]{%
				\hfill\makebox[10pt]{%
					\scriptsize\insertframenumber}}}}}
\newcommand\unnumbered{\setbeamertemplate{footline}{}}
\defbeamertemplate*{footline}{my infolines theme}
{
	\leavevmode%
	\hbox{%
		\begin{beamercolorbox}[wd=.33\paperwidth,ht=2.25ex,dp=1ex,center]{author in head/foot}%
			\usebeamerfont{author in head/foot}\insertshortauthor
		\end{beamercolorbox}%
		\begin{beamercolorbox}[wd=.34\paperwidth,ht=2.25ex,dp=1ex,center]{title in head/foot}%
			\usebeamerfont{title in head/foot}\insertshortinstitute
		\end{beamercolorbox}%
		\begin{beamercolorbox}[wd=.33\paperwidth,ht=2.25ex,dp=1ex,right]{date in head/foot}%
			\usebeamerfont{date in head/foot}\insertshortdate{}\hspace*{2em}
			\insertframenumber{} / \inserttotalframenumber\hspace*{2ex}
	\end{beamercolorbox}}%
	\vskip0pt%
}
%% --- Baskar's hack complete. Now give the following for each slide to be uncounted:
%% {\unnumbered\addtocounter{framenumber}{-1}


\AtBeginSection[]
{
	\begin{frame}\tiny
		\frametitle{Mnemonics}
		\unnumbered\addtocounter{framenumber}{-1}
		\tableofcontents[currentsection]
	\end{frame}
}
\begin{document}\definecolor{bbac35}{HTML}{bbac35}
	{
	\usebackgroundtemplate{\includegraphics[width=\paperwidth]{bitsgoatopslide.png}}
	\begin{frame}[plain]\frametitle{\color{red}Top Slide}% Make this empty in the final
 
 \vspace{.07\textheight}
 
 \begin{flushright}
 \colorbox{bbac35}{\textcolor{blue!100}{\huge\sc Presentation Title}}
 
 \colorbox{bbac35}{\textcolor{blue!100}{\large Subheading: Change the colors, sizes, font family etc.}}

 \colorbox{bbac35}{\textcolor{blue!100}{\Large \bbfamily Author Name}}

 \colorbox{bbac35}{\textcolor{blue!90}{\large \sc Affiliation Department etc.}}
 \end{flushright}
 	% This is only to get the image span the whole frame.
	\end{frame}
}
\setbeamertemplate{footline}
{
	{\includegraphics[width=0.5\paperwidth,height=1.5mm]{wilprezcoline.png}}
	~\raisebox{1.5ex}[0pt]{\insertsection\ \ \insertframenumber/\inserttotalframenumber}
}
\setbeamertemplate{headline}{.\hfill\includegraphics[scale=0.5]{wilprez3col.png}}
\begin{frame}{Welcome, Bienvenido, Bem-vinda, \hindifont{येवकार!}}

The Devanagari Konkani greeting in the frame title requires \textsc{Xe}\LaTeX\ to compile. Just \LaTeX\ wouldn't suffice.
\pause
\begin{center}
\includegraphics[height=1in]{airindiamaharaja.png}
\end{center}
\end{frame}
\frame{\tableofcontents}
\section{What is xyz}
\begin{frame}{xyz is ...}
   \begin{enumerate}[<+->]
       \item xyz is important
       \item xyz is useful
       \item xyz is new
   \end{enumerate}
\end{frame}
\note[itemize]{\item abc\item pqr}
\begin{frame}{More on xyz}
\begin{center}
 \includegraphics[width=.7\textwidth]{swagatsaral.png}

 This doesn't require \textsc{Xe}\LaTeX.
\end{center}
\end{frame}
\note{Discuss abc with xyz.}
\section{Bibliography}
\begin{frame}{Some Literature Survey} For the first citation, one needs a \texttt{\textbackslash citet*} to lay out a full author list, as \citet*{annie1} demonstrate, and this is not different for citing documents that have less than three authors as in \citet{shildesthesis}. One need not worry about citing several documents together, as in \citet{madhsardthesis,mcderm} or \citep{sinhala1,sinhala2} or \citealp{sampson1,sampson2}. One can also choose to cite Wirth\citeyearpar{wirth} with the year given differently.
    
\end{frame}
\begin{frame}[allowframebreaks]
	\frametitle{References}
	\bibliographystyle{plainnat}
	\bibliography{slides}
\end{frame}

\end{document}
